\paragraph{压力缺口的乘法射线势、escort 单比特实验与 Hellinger--Chernoff 母核}
在整数压力轴的冻结斜率恢复、escort--Rényi 熵率的压力差公式、末位统计的渐近充分化、escort 熵常数的闭式化以及 Fisher--Rao 几何已经确立之后，还可继续压缩出下列六条彼此闭合的后果。它们分别把乘法细化中的压力损失写成一条精确射线势，把可除偏序上的斜率单调与整倍数冻结传播压缩为纯序结构命题，把有界温度 escort 实验族在 Blackwell 意义下指数塌缩为单比特 Bernoulli 家族，并把全阶 Rényi 常数层、Shannon 温度耗散以及二温度判别几何统一到同一个 logistic 母核之下。

\begin{theorem}[压力缺口是乘法熵亏损的精确射线势]\label{thm:xi-time-part65h-pressure-gap-ray-potential}
沿用定理 \ref{thm:xi-time-part65c-escort-renyi-two-parameter-freezing-line} 与定理 \ref{thm:xi-time-part57b-pressure-spectrum-slope-intercept-recovery} 的记号。设整数压力极限
\[
P_a:=\lim_{m\to\infty}\frac1m\log S_a(m)
\qquad (a\ge 1)
\]
都存在，并记
\[
\alpha_*:=\lim_{q\to\infty}\frac{P_q}{q}.
\]
对每个整数 \(a\ge 1\) 记 escort 分布
\[
\pi_{m,a}(x):=\frac{d_m(x)^a}{S_a(m)}.
\]
再对整数 \(a\ge 1\)、\(b\ge 2\) 定义
\[
\delta(a):=\frac{P_a}{a}-\alpha_*,
\qquad
H_{\mathrm{loss}}(a,b):=bP_a-P_{ab}.
\]
再记 escort--Rényi 熵率
\[
h_b^{\mathrm{esc}}(a):=
\lim_{m\to\infty}\frac1m H_b(\pi_{m,a}).
\]
其存在性由定理 \ref{thm:xi-time-part65c-escort-renyi-two-parameter-freezing-line} 保证。则对任意整数 \(a\ge 1\)、\(b\ge 2\) 都有严格恒等式
\[
\delta(a)-\delta(ab)=\frac{H_{\mathrm{loss}}(a,b)}{ab}.
\]
并且
\[
\delta(a)-\delta(ab)=\frac{b-1}{ab}\,h_b^{\mathrm{esc}}(a)\ge 0.
\]
更进一步，沿任意乘法射线都有绝对收敛展开
\[
\delta(a)=\sum_{k=0}^{\infty}\frac{H_{\mathrm{loss}}(ab^k,b)}{ab^{k+1}}.
\]
\leanverified{paper\_xi\_time\_part65h\_pressure\_gap\_ray\_potential}
\end{theorem}
\begin{proof}
第一式只是代数恒等式：
\[
\delta(a)-\delta(ab)
=
\left(\frac{P_a}{a}-\alpha_*\right)
-\left(\frac{P_{ab}}{ab}-\alpha_*\right)
=
\frac{bP_a-P_{ab}}{ab}
=
\frac{H_{\mathrm{loss}}(a,b)}{ab}.
\]

另一方面，由定理 \ref{thm:xi-time-part65c-escort-renyi-two-parameter-freezing-line}，
\[
h_b^{\mathrm{esc}}(a)
=
\frac{bP_a-P_{ab}}{b-1}
=
\frac{H_{\mathrm{loss}}(a,b)}{b-1}.
\]
代回即得
\[
\delta(a)-\delta(ab)=\frac{b-1}{ab}\,h_b^{\mathrm{esc}}(a)\ge 0,
\]
其中非负性来自 Rényi 熵率本身非负。

最后，对任意 \(N\ge 1\) 作有限望远镜求和：
\[
\delta(a)-\delta(ab^N)
=
\sum_{k=0}^{N-1}\bigl(\delta(ab^k)-\delta(ab^{k+1})\bigr)
=
\sum_{k=0}^{N-1}\frac{H_{\mathrm{loss}}(ab^k,b)}{ab^{k+1}}.
\]
由定理 \ref{thm:xi-time-part57b-pressure-spectrum-slope-intercept-recovery}，
\[
\frac{P_q}{q}\longrightarrow \alpha_*
\qquad (q\to\infty),
\]
故当 \(N\to\infty\) 时，
\[
\delta(ab^N)=\frac{P_{ab^N}}{ab^N}-\alpha_*\longrightarrow 0.
\]
于是得到
\[
\delta(a)=\sum_{k=0}^{\infty}\frac{H_{\mathrm{loss}}(ab^k,b)}{ab^{k+1}}.
\]
又由于各项都可写成
\[
\frac{H_{\mathrm{loss}}(ab^k,b)}{ab^{k+1}}
=
\frac{b-1}{ab^{k+1}}\,h_b^{\mathrm{esc}}(ab^k)\ge 0,
\]
该级数因而绝对收敛。证毕。
\end{proof}

\begin{corollary}[可除偏序上的斜率单调、几何逼近与整倍数冻结传播]\label{cor:xi-time-part65h-divisibility-slope-freezing-propagation}
沿用定理 \ref{thm:xi-time-part65h-pressure-gap-ray-potential} 的记号。则成立：
\begin{enumerate}
  \item 若 \(a\mid c\)，则
  \[
  \frac{P_c}{c}\le \frac{P_a}{a}.
  \]
  \item 对任意整数 \(a\ge 1\)、\(b\ge 2\)、\(N\ge 0\)，都有
  \[
  0\le \frac{P_{ab^N}}{ab^N}-\alpha_*
  \le \frac{\log\varphi}{ab^N}.
  \]
  \item 若对某个整数 \(a_0\ge 1\) 恰有
  \[
  P_{a_0}=a_0\alpha_*,
  \]
  则对每个 \(n\ge 1\) 都有
  \[
  P_{a_0n}=a_0n\,\alpha_*.
  \]
\end{enumerate}
\end{corollary}
\begin{proof}
若 \(c=a\)，结论平凡。若 \(c>a\) 且 \(a\mid c\)，写 \(c=ab\) 且 \(b\ge 2\)。由定理 \ref{thm:xi-time-part65h-pressure-gap-ray-potential}，
\[
\frac{P_a}{a}-\frac{P_c}{c}
=
\delta(a)-\delta(ab)
=
\frac{b-1}{ab}\,h_b^{\mathrm{esc}}(a)\ge 0,
\]
故第 \(1\) 条成立。

对第 \(2\) 条，令 \(n:=ab^N\)。由定义，
\[
M_m^n\le S_n(m)\le |X_m|\,M_m^n.
\]
两边取对数后除以 \(mn\)，再令 \(m\to\infty\)，得到
\[
\alpha_*\le \frac{P_n}{n}\le \alpha_*+\frac{1}{n}\lim_{m\to\infty}\frac1m\log|X_m|.
\]
而 \(|X_m|=F_{m+2}\)，故
\[
\lim_{m\to\infty}\frac1m\log|X_m|=\log\varphi.
\]
于是
\[
0\le \frac{P_n}{n}-\alpha_*\le \frac{\log\varphi}{n}.
\]
代回 \(n=ab^N\) 即得所示估计。

最后，若 \(P_{a_0}=a_0\alpha_*\)，则对任意 \(n\ge 1\)，由第 \(1\) 条与第 \(2\) 条分别得到
\[
\alpha_*\le \frac{P_{a_0n}}{a_0n}\le \frac{P_{a_0}}{a_0}=\alpha_*.
\]
故必有
\[
\frac{P_{a_0n}}{a_0n}=\alpha_*,
\]
即
\[
P_{a_0n}=a_0n\,\alpha_*.
\]
证毕。
\end{proof}

\begin{theorem}[有界温度 escort 实验族与单比特 Bernoulli 实验的 Blackwell 指数等价]\label{thm:xi-time-part65h-escort-blackwell-onebit-equivalence}
沿用定理 \ref{thm:killo-fold-bin-escort-onebit-fisher} 的记号。固定 \(Q<\infty\)，定义统计实验族
\[
\mathcal E_m^{(Q)}:=\{\pi_{m,q}:0\le q\le Q\},
\]
以及极限单比特实验族
\[
\mathcal B^{(Q)}:=\{\beta_q:0\le q\le Q\},
\qquad
\beta_q:=\mathrm{Bernoulli}(\theta(q)),
\qquad
\theta(q):=\frac{1}{1+\varphi^{q+1}}.
\]
再定义
\[
\Delta_Q(m):=
\inf_{K,L}\sup_{0\le q\le Q}
\max\Bigl\{
D_{\mathrm{TV}}(K_{\#}\pi_{m,q},\beta_q),\
D_{\mathrm{TV}}(L_{\#}\beta_q,\pi_{m,q})
\Bigr\},
\]
其中 \(K:X_m\rightsquigarrow\{0,1\}\) 与 \(L:\{0,1\}\rightsquigarrow X_m\) 取遍全部 Markov 核。则存在常数 \(C_Q>0\) 使得
\[
\Delta_Q(m)\le C_Q\Bigl(\frac{\varphi}{2}\Bigr)^m.
\]
因此，对整个有界温度区间而言，escort 统计实验在 Blackwell 意义下以指数速度塌缩为由末位比特承载的一维 logistic Bernoulli 家族。
\end{theorem}
\begin{proof}
记
\[
A_i:=\{w\in X_m:w_m=i\},
\qquad
U_{m,i}:=\mathrm{Unif}(A_i)
\qquad (i=0,1).
\]
定义前向核
\[
K_m(w):=w_m\in\{0,1\},
\]
以及反向核
\[
L_m(i,\cdot):=U_{m,i}.
\]
由定理 \ref{thm:killo-fold-bin-escort-onebit-fisher}，存在两态 Gibbs 近似
\[
\nu_{m,q}(w):=\frac{u_m(w)\varphi^{-qw_m}}{Z_{m,q}},
\qquad
Z_{m,q}:=\frac{F_{m+1}+\varphi^{-q}F_m}{F_{m+2}},
\]
并且
\[
\sup_{0\le q\le Q}\|\pi_{m,q}-\nu_{m,q}\|_{\mathrm{TV}}
=
O_Q\!\Bigl(\Bigl(\frac{\varphi}{2}\Bigr)^m\Bigr).
\]
由于 \(\nu_{m,q}\) 在每个块 \(A_i\) 上严格均匀，可把它写成
\[
\nu_{m,q}=(1-\theta_m(q))U_{m,0}+\theta_m(q)U_{m,1},
\]
其中
\[
\theta_m(q):=\nu_{m,q}(A_1)
=
\frac{\varphi^{-q}F_m}{F_{m+1}+\varphi^{-q}F_m}.
\]

先看前向方向。由 \(K_m\) 的定义，
\[
K_{m\#}\pi_{m,q}=\mathrm{Bernoulli}(\theta_m(q)).
\]
而
\[
\theta_m(q)=\frac{\varphi^{-q}}{F_{m+1}/F_m+\varphi^{-q}},
\qquad
\theta(q)=\frac{\varphi^{-q}}{\varphi+\varphi^{-q}}.
\]
又因为
\[
\frac{F_{m+1}}{F_m}=\varphi+O(\varphi^{-2m}),
\]
并且 \(q\in[0,Q]\) 落在紧区间上，所以
\[
\sup_{0\le q\le Q}|\theta_m(q)-\theta(q)|
=
O_Q(\varphi^{-2m})
=
O_Q\!\Bigl(\Bigl(\frac{\varphi}{2}\Bigr)^m\Bigr).
\]
因此
\[
\sup_{0\le q\le Q}
D_{\mathrm{TV}}(K_{m\#}\pi_{m,q},\beta_q)
=
O_Q\!\Bigl(\Bigl(\frac{\varphi}{2}\Bigr)^m\Bigr).
\]

再看反向方向。由构造，
\[
L_{m\#}\beta_q=(1-\theta(q))U_{m,0}+\theta(q)U_{m,1}.
\]
而与 \(\nu_{m,q}\) 相比，两者仅块权重不同，故
\[
\|L_{m\#}\beta_q-\nu_{m,q}\|_{\mathrm{TV}}
=
|\theta(q)-\theta_m(q)|
=
O_Q\!\Bigl(\Bigl(\frac{\varphi}{2}\Bigr)^m\Bigr).
\]
再与 \(\pi_{m,q}\) 到 \(\nu_{m,q}\) 的估计合并，得到
\[
\sup_{0\le q\le Q}
D_{\mathrm{TV}}(L_{m\#}\beta_q,\pi_{m,q})
=
O_Q\!\Bigl(\Bigl(\frac{\varphi}{2}\Bigr)^m\Bigr).
\]
对 \(\Delta_Q(m)\) 取上述这对核 \((K_m,L_m)\) 作为候选，即得所示上界。证毕。
\end{proof}

\begin{theorem}[bin-fold escort 的全阶 Rényi 熵在常数层具有显式温度剖面]\label{thm:xi-time-part65h-escort-renyi-o1-temperature-profile}
沿用定理 \ref{thm:killo-fold-bin-escort-onebit-fisher} 的记号。对固定 \(q\ge 0\) 与 \(\alpha>0\)，定义 escort Rényi 熵
\[
H_{\alpha}(\pi_{m,q})
:=
\begin{cases}
\dfrac{1}{1-\alpha}\log\sum_{w\in X_m}\pi_{m,q}(w)^\alpha,
& \alpha\neq 1,\\[2mm]
-\sum_{w\in X_m}\pi_{m,q}(w)\log\pi_{m,q}(w),
& \alpha=1.
\end{cases}
\]
则当 \(m\to\infty\) 时，若 \(\alpha\neq 1\)，有
\[
H_{\alpha}(\pi_{m,q})
=
m\log\varphi
-\frac12\log 5
+
\frac{\alpha\log(\varphi+\varphi^{-q})-\log(\varphi+\varphi^{-\alpha q})}{\alpha-1}
+O\!\Bigl(\Bigl(\frac{\varphi}{2}\Bigr)^m\Bigr),
\]
而当 \(\alpha=1\) 时，有
\[
H(\pi_{m,q})
=
m\log\varphi
-\frac12\log 5
+\log(\varphi+\varphi^{-q})
+\frac{q\log\varphi}{1+\varphi^{q+1}}
+O\!\Bigl(\Bigl(\frac{\varphi}{2}\Bigr)^m\Bigr).
\]
因此，escort 温度族不只在熵率层面塌缩到 \(\log\varphi\)；其全部 \(O(1)\) 剩余结构也由 \(q\) 的显式闭式函数完全控制。
\end{theorem}
\begin{proof}
先设 \(\alpha\neq 1\)。由定义，
\[
\sum_{w\in X_m}\pi_{m,q}(w)^\alpha
=
\frac{S_{\alpha q}^{\mathrm{bin}}(m)}{(S_q^{\mathrm{bin}}(m))^\alpha}.
\]
而命题 \ref{prop:fold-bin-power-sum-asymptotic} 给出，对每个固定 \(t\ge 0\)，
\[
S_t^{\mathrm{bin}}(m)
=
\Bigl(\frac{2}{\varphi}\Bigr)^{tm}
\bigl(F_{m+1}+\varphi^{-t}F_m\bigr)
\Bigl(1+O\!\Bigl(\Bigl(\frac{\varphi}{2}\Bigr)^m\Bigr)\Bigr).
\]
将 \(t=q\) 与 \(t=\alpha q\) 代入，即得
\[
\log\sum_{w\in X_m}\pi_{m,q}(w)^\alpha
=
\log\bigl(F_{m+1}+\varphi^{-\alpha q}F_m\bigr)
-\alpha\log\bigl(F_{m+1}+\varphi^{-q}F_m\bigr)
+O\!\Bigl(\Bigl(\frac{\varphi}{2}\Bigr)^m\Bigr).
\]
再由 Binet 公式，
\[
F_{m+1}+\varphi^{-u}F_m
=
\frac{\varphi^m}{\sqrt5}\bigl(\varphi+\varphi^{-u}\bigr)+O(\varphi^{-m}),
\]
故
\[
\log\bigl(F_{m+1}+\varphi^{-u}F_m\bigr)
=
m\log\varphi-\frac12\log5+\log(\varphi+\varphi^{-u})
+O(\varphi^{-2m}).
\]
代回并除以 \(1-\alpha\)，便得到所示 Rényi 公式。

再看 Shannon 情形。由定理 \ref{thm:xi-fold-escort-renyi-shannon-constants} 的证明所用恒等式，
\[
H(\pi_{m,q})
=
\log S_q^{\mathrm{bin}}(m)-q\,\kappa_m^{\mathrm{bin}}(q),
\]
其中
\[
\kappa_m^{\mathrm{bin}}(q):=
\sum_{w\in X_m}\pi_{m,q}(w)\log d_m^{\mathrm{bin}}(w).
\]
由推论 \ref{cor:xi-time-part9o-escort-fiberinfo-constantterm-inversion}，
\[
\kappa_m^{\mathrm{bin}}(q)
=
m\log\!\Bigl(\frac{2}{\varphi}\Bigr)
-\frac{\log\varphi}{1+\varphi^{q+1}}
+O\!\Bigl(\Bigl(\frac{\varphi}{2}\Bigr)^m\Bigr).
\]
另一方面，把 \(t=q\) 代入上述幂和展开并取对数，得到
\[
\log S_q^{\mathrm{bin}}(m)
=
qm\log\!\Bigl(\frac{2}{\varphi}\Bigr)
+m\log\varphi
-\frac12\log5
+\log(\varphi+\varphi^{-q})
+O\!\Bigl(\Bigl(\frac{\varphi}{2}\Bigr)^m\Bigr).
\]
两式合并即得
\[
H(\pi_{m,q})
=
m\log\varphi
-\frac12\log5
+\log(\varphi+\varphi^{-q})
+\frac{q\log\varphi}{1+\varphi^{q+1}}
+O\!\Bigl(\Bigl(\frac{\varphi}{2}\Bigr)^m\Bigr).
\]
证毕。
\end{proof}

\begin{theorem}[escort Shannon 常数满足精确 de Bruijn--Fisher 身份]\label{thm:xi-time-part65h-escort-shannon-debruijn-fisher}
沿用定理 \ref{thm:xi-time-part65h-escort-renyi-o1-temperature-profile} 的记号。定义 Shannon 常数层
\[
\mathfrak h(q):=
\lim_{m\to\infty}\bigl(H(\pi_{m,q})-m\log\varphi\bigr).
\]
则有闭式
\[
\mathfrak h(q)
=
-\frac12\log5+\log(\varphi+\varphi^{-q})
+\frac{q\log\varphi}{1+\varphi^{q+1}}.
\]
进一步，记
\[
\theta(q):=\frac{1}{1+\varphi^{q+1}},
\qquad
\mathcal I(q):=(\log\varphi)^2\theta(q)\bigl(1-\theta(q)\bigr).
\]
则成立精确微分恒等式
\[
\mathfrak h'(q)=-q\,\mathcal I(q),
\]
以及积分表示
\[
\mathfrak h(q)=\mathfrak h(0)-\int_0^q t\,\mathcal I(t)\,\dd t.
\]
其中 \(\mathcal I(q)\) 正是结论 \ref{con:xi-time-part9o-escort-fisher-rao-closed} 给出的极限 Fisher 信息密度。
\end{theorem}
\begin{proof}
第一条公式正是定理 \ref{thm:xi-time-part65h-escort-renyi-o1-temperature-profile} 在 \(\alpha=1\) 时的常数项。

令
\[
A(q):=\varphi+\varphi^{-q}.
\]
则
\[
\theta(q)=\frac{\varphi^{-q}}{A(q)}=\frac{1}{1+\varphi^{q+1}},
\]
并且
\[
\mathfrak h(q)
=
-\frac12\log5+\log A(q)+q\,\theta(q)\log\varphi.
\]
直接求导得
\[
\frac{A'(q)}{A(q)}
=
-\frac{\varphi^{-q}\log\varphi}{A(q)}
=
-\theta(q)\log\varphi,
\]
而
\[
\theta'(q)
=
-(\log\varphi)\theta(q)\bigl(1-\theta(q)\bigr).
\]
因此
\[
\mathfrak h'(q)
=
-\theta(q)\log\varphi+\theta(q)\log\varphi+q\,\theta'(q)\log\varphi
=
-q(\log\varphi)^2\theta(q)\bigl(1-\theta(q)\bigr)
=
-q\,\mathcal I(q).
\]
最后对 \(q\) 从 \(0\) 到任意给定值积分，即得
\[
\mathfrak h(q)=\mathfrak h(0)-\int_0^q t\,\mathcal I(t)\,\dd t.
\]
证毕。
\end{proof}

\begin{theorem}[escort 温度间的 Hellinger--Chernoff 几何由闭式母核完全支配]\label{thm:xi-time-part65h-escort-hellinger-chernoff-kernel}
沿用定理 \ref{thm:killo-fold-bin-escort-onebit-fisher} 的记号。对任意 \(p,q\ge 0\) 与 \(s\in[0,1]\)，定义混合亲和量
\[
\mathcal M_s^{(m)}(p,q):=
\sum_{w\in X_m}\pi_{m,p}(w)^s\pi_{m,q}(w)^{1-s}.
\]
则极限
\[
\mathcal M_s(p,q):=\lim_{m\to\infty}\mathcal M_s^{(m)}(p,q)
\]
存在，并满足精确闭式
\[
\mathcal M_s(p,q)
=
\frac{\varphi+\varphi^{-(sp+(1-s)q)}}
{(\varphi+\varphi^{-p})^s(\varphi+\varphi^{-q})^{1-s}}.
\]
特别地，Hellinger 亲和量
\[
\mathcal H(p,q):=\mathcal M_{1/2}(p,q)
\]
满足
\[
\mathcal H(p,q)
=
\frac{\varphi+\varphi^{-(p+q)/2}}
{\sqrt{(\varphi+\varphi^{-p})(\varphi+\varphi^{-q})}},
\]
而 Chernoff 母核可写为
\[
\mathcal C(p,q):=\inf_{0\le s\le 1}\mathcal M_s(p,q).
\]
因此，两温度 escort 之间的全部 Hellinger--Chernoff 判别几何都由单参数显式核
\[
s\longmapsto \mathcal M_s(p,q)
\]
完全控制。
\end{theorem}
\begin{proof}
对任意 \(s\in[0,1]\)，由 escort 分布的定义，
\[
\pi_{m,p}(w)^s\pi_{m,q}(w)^{1-s}
=
\frac{d_m^{\mathrm{bin}}(w)^{sp+(1-s)q}}
{\bigl(S_p^{\mathrm{bin}}(m)\bigr)^s\bigl(S_q^{\mathrm{bin}}(m)\bigr)^{1-s}}.
\]
故有严格恒等式
\[
\mathcal M_s^{(m)}(p,q)
=
\frac{S_{sp+(1-s)q}^{\mathrm{bin}}(m)}
{\bigl(S_p^{\mathrm{bin}}(m)\bigr)^s\bigl(S_q^{\mathrm{bin}}(m)\bigr)^{1-s}}.
\]
对分子与分母分别应用命题 \ref{prop:fold-bin-power-sum-asymptotic}，
\[
S_t^{\mathrm{bin}}(m)
=
\Bigl(\frac{2}{\varphi}\Bigr)^{tm}
\bigl(F_{m+1}+\varphi^{-t}F_m\bigr)
\Bigl(1+O\!\Bigl(\Bigl(\frac{\varphi}{2}\Bigr)^m\Bigr)\Bigr)
\qquad (t\ge 0).
\]
其中指数因子完全抵消。
于是
\[
\mathcal M_s^{(m)}(p,q)
=
\frac{F_{m+1}+\varphi^{-(sp+(1-s)q)}F_m}
{\bigl(F_{m+1}+\varphi^{-p}F_m\bigr)^s
\bigl(F_{m+1}+\varphi^{-q}F_m\bigr)^{1-s}}
\Bigl(1+O\!\Bigl(\Bigl(\frac{\varphi}{2}\Bigr)^m\Bigr)\Bigr).
\]
再由 Binet 公式
\[
F_{m+1}+\varphi^{-u}F_m
=
\frac{\varphi^m}{\sqrt5}\bigl(\varphi+\varphi^{-u}\bigr)+O(\varphi^{-m}),
\]
便得到
\[
\mathcal M_s^{(m)}(p,q)
=
\frac{\varphi+\varphi^{-(sp+(1-s)q)}}
{(\varphi+\varphi^{-p})^s(\varphi+\varphi^{-q})^{1-s}}
+O\!\Bigl(\Bigl(\frac{\varphi}{2}\Bigr)^m\Bigr).
\]
令 \(m\to\infty\) 即得所示极限核。

取 \(s=\tfrac12\) 立即得到 Hellinger 亲和量公式；而 Chernoff 母核不过是同一闭式核在区间 \([0,1]\) 上的下确界。证毕。
\end{proof}

由此，escort 主线中的整数压力差、末位单比特充分化、Rényi 常数层、Fisher 温度耗散与 Hellinger--Chernoff 判别核便被进一步压缩到同一条闭合链上：乘法偏序一侧的全部亏损由射线势 \(\delta\) 精确编码，而统计实验一侧的全部有限温度几何则完全退化为由 \(\theta(q)=1/(1+\varphi^{q+1})\) 所驱动的一维 logistic 母模型。
